\section{Introduction}\label{sec:introduction}


Let a smooth family be completed to a morphism of projective pairs.
How many logarithmic pluricanonical forms can its total space have?
Subadditivity gives a lower bound from the base and a general fiber.
We prove the complementary upper bound when the total space and all
boundary strata are smooth over the open base. The boundary may have
horizontal components, and the conclusion includes the case in which
the base has no logarithmic pluricanonical forms at all.

Throughout, varieties are over $\C$. Let $f:X\to Y$ be a surjective
morphism with connected fibers between smooth connected projective
varieties. Let $E$ and $D$ be reduced simple-normal-crossing divisors
on $X$ and $Y$, respectively, with
\begin{equation}\label{eq:support-inclusion}
 \Supp(f^*D)\subseteq\Supp(E).
\end{equation}
Either divisor may be zero. Set $V=Y\setminus\Supp D$. We assume that
$X$ and every irreducible component of every nonempty intersection of
components of $E$, restricted to $f^{-1}(V)$, are smooth over $V$.
We call this \emph{stratum smoothness over $V$}. For a very general
$y\in V$, set $F=X_y$ and $E_F=E|_F$. Write
\[
 L_X=K_X+E,\qquad L_Y=K_Y+D,\qquad L_F=K_F+E_F.
\]
We give a point Kodaira dimension zero and use
$(-\infty)+a=-\infty$, also when $a=-\infty$.

\begin{theorem}[Reverse logarithmic Kodaira inequality]\label{thm:upper-bound}
Under these hypotheses,
\begin{equation}\label{eq:upper-bound}
 \kappa(X,L_X)\leq\kappa(Y,L_Y)+\kappa(F,L_F).
\end{equation}
If either term on the right is $-\infty$, then
$H^0(X,mL_X)=0$ for every $m>0$.
\end{theorem}

There is no abundance, semiampleness, good-minimal-model, or general-type
hypothesis. Condition \eqref{eq:support-inclusion} is a condition on supports;
the multiplicities in $f^*D$ remain in the proof. No smoothness is
assumed over $D$.

The additional input for equality is the logarithmic subadditivity
theorem of the companion article \cite[Corollary~\FoundationLogNumber]{Foundation}.
It applies to any connected-fiber surjective morphism of smooth
projective reduced-SNC pairs satisfying \eqref{eq:support-inclusion}, without
the stratum-smoothness assumption. Its fiber is the same very general
pair $(F,E_F)$, and it gives the opposite inequality. Therefore:

\begin{corollary}[Logarithmic additivity]\label{cor:additivity}
For the pairs in Theorem~\ref{thm:upper-bound},
\begin{equation}\label{eq:additivity}
 \kappa(X,K_X+E)=\kappa(Y,K_Y+D)+\kappa(F,K_F+E_F).
\end{equation}
\end{corollary}

This gives a positive resolution of Popa's logarithmic additivity
conjecture in its projective
reduced-SNC, stratum-smooth formulation
\cite[Conjecture 3.9 and its defining footnote]{PopaSurvey}.
The proof of the upper inequality is independent of logarithmic
subadditivity; the companion theorem is used only for this corollary.

\subsection{The section-construction result}

The difficult case is a base with logarithmic Kodaira dimension zero.
The space of logarithmic pluricanonical sections then has dimension
at most one in each degree. Choose a very general point outside the zero divisors of all
these sections. If a total-space pluriform vanished on its fiber, we
would like to produce a base pluriform vanishing at that point. This
would be a contradiction, and restriction to the fiber would be
injective in every degree.

The following result carries out that construction. It also treats a
base of logarithmic Kodaira dimension $-\infty$.

\begin{proposition}\label{prop:section-construction}
Let $f,E,D$ satisfy the hypotheses of Theorem~\ref{thm:upper-bound}, and
let $0\ne s\in H^0(X,mL_X)$ for some $m>0$. Then
$\kappa(Y,L_Y)\geq0$. If $\kappa(Y,L_Y)=0$ and $s$ vanishes identically
on $X_y$ for some $y\in V$, there are an integer $a>0$ and a nonzero
$\tau\in H^0(Y,aL_Y)$ with $\tau(y)=0$.
\end{proposition}

The point $y$ in this statement belongs to the original open base
$V$. It need not belong to a smaller open set chosen in constructing
a cyclic cover from $s$. This uniformity is what allows the restriction
argument to handle all total-space sections at one very general fiber.
In particular, if $\kappa(Y,L_Y)=0$, restriction to one very general
fiber is injective in every positive degree, so
\[
 h^0(X,mL_X)\leq h^0(F,mL_F)
 \qquad\text{for every positive integer }m.
\]
\Cref{sec:completion} proves \Cref{thm:upper-bound} from this proposition,
including both negative-infinity cases.

\subsection{Two outcomes of the Hodge construction}

A pluriform on the total space first produces a cohomology class:
take a root of the pluriform on a finite cover, then project its class
to a nonzero pure weight quotient. This gives a vector in the bundle
of highest Hodge classes. The chosen vector need not span that bundle.
Thus positivity of the determinant of the whole bundle does not yet
control the coordinates of this particular vector, or retain its zero
on a chosen fiber.

We keep track of these coordinates in the logarithmic differential
frames of the original pair $(Y,D)$. The root cover may acquire extra
singular fibers inside $V$, but its new discriminant must not enlarge
the poles allowed on the base. A local calculation shows both that
no new poles occur and that a zero on any original smooth fiber is
retained. Homogeneous tensor operations preserve these properties
while making the cohomological data descend to a projective parameter
space $S$. This space records the varying Hodge structure and the
additional vector directions that its period map alone can miss.
There are two possibilities.

If $S$ is a point, the resulting cohomology bundle is constant.
Its vector has ordinary scalar coordinates, at least one of which
is nonzero. The coefficient calculation makes each coordinate an
actual base pluriform. The same calculation preserves the prescribed
zero. This proves base nonvanishing, excluding a base of logarithmic Kodaira
dimension $-\infty$. The preserved zero gives the additional conclusion
needed when its logarithmic Kodaira dimension is zero.

If $S$ has positive dimension, restrict first to the generic fiber
of the map from the base to $S$. The coefficient system is constant
there, so a nonzero scalar coordinate permits a relative Iitaka
fibration. It produces a divisor on an intermediate variety with the
same sufficiently divisible section spaces as $K_Y+D$. Period
positivity makes this divisor big after adding a pullback from $S$.
The remaining obstacle is to remove that added divisor. We do so by
comparing the iterated Higgs images of the chosen vector with the
period directions and applying a numerical subtraction argument.
The divisor itself is then big, which forces
$\kappa(Y,K_Y+D)>0$. Thus this alternative cannot occur when the
base has logarithmic Kodaira dimension zero or $-\infty$.

\subsection{Earlier results and the proof route}

Iitaka's subadditivity problem asks whether
the Kodaira dimension of an algebraic fiber space is at least the sum
of those of its base and general fiber. Viehweg proved the
relative-dimension-one case \cite[Introduction]{Viehweg1977}.
Logarithmic pluriforms bring open varieties into the same birational
framework \cite[Introduction]{IitakaLogForms1976}.

For smooth families, the stronger question is whether the lower bound
becomes equality. Popa's logarithmic formulation specifies smoothness
of the total space and every boundary stratum over the open base
\cite[Conjecture 3.9 and footnote 5]{PopaSurvey}. The reverse inequality
constrains the extra pluricanonical growth that variation might
otherwise produce. Its negative-base case makes the distinction from
subadditivity especially clear: the lower bound is then automatic,
whereas the upper bound asserts vanishing of every positive
pluricanonical space on the total pair.

Popa--Schnell proved the upper inequality for smooth projective
families under the Campana--Peternell conjecture and obtained a
degreewise plurigenus bound when the base has rationally trivial
canonical bundle \cite[Theorems C and H]{PopaSchnellSmooth}.
Their reduction through the base Iitaka fibration also organizes our
argument. The section-construction proposition adds the specific
conclusion needed here: an actual base pluriform, with a prescribed
zero whenever the original section vanishes on a chosen smooth fiber.

Fujino--Fujisawa's logarithmic results already allow relatively
normal-crossing horizontal boundary. Their upper inequality for a base of nonnegative logarithmic Kodaira
dimension assumes generalized abundance for sufficiently general fibers of the base's
logarithmic Iitaka fibration \cite[Theorem 1.12]{FF}; their Hodge-theoretic construction also
descends pseudoeffectivity to the base
\cite[Theorems 1.4 and 4.5]{FF}.
Park proves the logarithmic upper inequality, in the reduced
inverse-image boundary setting and for a base of nonnegative
logarithmic Kodaira dimension, assuming good minimal models for the
very general fibers of that base Iitaka fibration
\cite[Theorem 1.12]{Park}. These hypotheses concern fibers of the
\emph{base's} Iitaka fibration, rather than the original fibers of $f$.

A different route is Campana's logarithmic additivity theorem for
smooth projective families over a quasiprojective base with semiample
canonical bundles on the original fibers
\cite[Theorem 1]{CampanaAdditivity}. Its proof uses birational
isotriviality and the core. The extension to original fibers with good
minimal models uses Taji's birational-isotriviality theorem over a
special base \cite[Theorem 1.1]{Taji}, as explained in
\cite[Remark 2(1)]{CampanaAdditivity}.
Theorem~\ref{thm:upper-bound} imposes neither of these minimal-model
conditions.

The differential construction has an earlier model in Kov\'acs's
iteration of cohomology maps from logarithmic cotangent sequences
\cite[Section 1, Lemma 1.1]{KovacsFineModuli}. The passage from a
pluriform to a cyclic cover and iterated Higgs maps belongs to the
Viehweg--Zuo method
\cite[Introduction, Lemma 1.1, and Section 3]{ViehwegZuoIsotriviality}.
We use its Hodge-theoretic development in Popa--Schnell and the
logarithmic construction of Fujino--Fujisawa, while retaining the
chosen vector and its original coefficient lattice. Further inputs
are full integral period-image algebraicity \cite{BBT}, orbifold
cotangent positivity \cite{CP}, and b-nefness of the moduli part of
an lc-trivial fibration \cite{FG}. For the last input we verify the
rank condition for an auxiliary subpair which may have negative
coefficients. The new task addressed by the section construction is
the passage from these positivity statements to actual base sections
with the stated control at every point of the original open base.

\paragraph{Organization.}
\Cref{sec:completion} deduces the upper inequality in all sign cases from
\Cref{prop:section-construction}. The rest of the paper proves that
proposition. \Cref{sec:logarithmic-input} constructs the geometric vector,
its coefficient lattice, and the zero test. \Cref{sec:marked-periods}
constructs the quotient that retains its periods and compact-factor
directions. \Cref{sec:boundary-rank} proves the two comparisons that remove
the boundary and control the chosen vector by its iterated Higgs images.
\Cref{sec:relative-iitaka} produces an intermediate base with exactly the
original logarithmic section spaces. \Cref{sec:nonvanishing} first proves
the constant-parameter case, then removes the period twist in the
positive-dimensional case and completes the section construction. Metric and
flat-connection estimates appear immediately before their first uses.
